\documentclass{situationist-preprint}

\title{Promises Without Reliance}
\msid{ETH-16-b}
\msfamily{ETH-16}
\msdiscipline{Normative and Applied Ethics}
\mstype{Consequence}
\mskeywords{promising; reliance; expectation view; unwanted promises; coerced promises}

\begin{document}
\maketitle

\begin{abstract}
ETH-16-a argues that promissory obligation arises from the creation of reasonable reliance and persists when the promisee forgets. A standard objection to reliance views concerns promises on which the promisee never relies: she does not believe the promise, or does not change her plans on account of it. We derive the verdicts of ETH-16-a for these cases. Where the promisee does not actually rely but is entitled to, the promise binds directly, since what the promisor creates is the entitlement and not the psychological state. Where no entitlement arises, because the promisee does not want the assurance, because the promise was extorted by her, or because there is no other party at all, the promise does not bind as a promise, though it may bind derivatively through other duties.
\end{abstract}

\section{Two senses of reliance}

ETH-16-a distinguishes \emph{reasonable reliance}, the promisee's entitlement, in virtue of what the promisor has done, to expect the performance and plan on it, from \emph{actual reliance}, the psychological state of expecting and planning (Definition~2.2 of ETH-16-a). Its reliance principle (Principle~P1 of ETH-16-a) grounds the obligation in the voluntary and intentional creation of reasonable reliance, where the promisor knows the promisee wants assurance; and Lemma~3.1 of ETH-16-a shows that every promise creates reasonable reliance in the promisee. A \enquote{promise without reliance} may therefore be one of two kinds: one that creates no actual reliance, and one in which the conditions under which the creation of reasonable reliance obligates are not met. The account treats them differently.

\section{No actual reliance}

\begin{proposition}\label{p:noactual}
A promise binds although the promisee never actually relies on it, provided the conditions of Principle~P1 of ETH-16-a are met.
\end{proposition}

\begin{proof}
By Lemma~3.1 and Theorem~3.2 of ETH-16-a, the obligation arises from the creation of reasonable reliance, which every promise creates. Actual reliance does not figure among the conditions of Principle~P1 of ETH-16-a. Its absence therefore does not prevent the obligation from arising.
\end{proof}

Two cases illustrate the proposition. \emph{The skeptical promisee}: a friend who has broken many promises promises to repay a loan by the end of the month, and the lender, from experience, does not believe him and makes her plans without counting on the money. \emph{The self-sufficient promisee}: a colleague promises to bring a projector to a meeting, and the speaker, who wants the assurance, brings her own anyway. In both cases the promisee is entitled to expect performance and to plan on it, whether or not she does; and in both, the promisor who fails to perform cannot excuse himself by pointing out that she did not count on him. Her disbelief or caution is a response to his unreliability or to the stakes, not a waiver of what he has given her. The proposition is the counterpart, at the moment of promising, of the persistence theorem of ETH-16-a (Theorem~3.3), which shows that the obligation survives the loss of actual reliance after it has formed.

\section{Where the reliance principle fails}

Principle~P1 of ETH-16-a obligates only where reasonable reliance is created voluntarily and intentionally, by a promisor who knows that a promisee distinct from him wants assurance. Three cases fail these conditions.

\paragraph{Unwanted promises.} A neighbor promises to visit every Sunday; the person to whom he makes the promise does not want the visits, and he knows it. Principle~P1 of ETH-16-a applies only where the promisor knows that the promisee wants assurance. Here she does not, and the promisor's act gives her nothing she has reason to count on in her own interest. No promissory obligation to her arises. If the promise was made in public, and others plan on the visits, any duty the neighbor has is owed to them under the general wrong of defeating expectations one has deliberately created \citep{Scanlon1998}, which binds only derivatively and only to the extent of their reliance.

\paragraph{Extorted promises.} A promise extracted from the promisor by the promisee's threat is not made voluntarily, and whatever expectation the promisee has she owes to her own coercion. The condition of voluntariness in Principle~P1 of ETH-16-a is not met, and the promise does not bind, a verdict that conventionalists from \citet[3.2.5]{Hume1739} onward have also accepted.

\paragraph{Promises to oneself.} A resolution to exercise daily, described as a promise to oneself, has no promisee distinct from the promisor. There is no promise in the sense of Definition~2.1 of ETH-16-a, which requires uptake by another, and ETH-16-a already classes uncommunicated resolutions as non-binding. Such promises bind not at all as promises; what binds is whatever reason the agent had for the resolution.

\begin{proposition}\label{p:noreasonable}
Where the conditions of Principle~P1 of ETH-16-a are not met, because the promisee does not want assurance, the act is not voluntary, or there is no promisee distinct from the promisor, no promissory obligation arises, though the agent may be bound derivatively by other duties to those who rely on the act.
\end{proposition}

\begin{proof}
By Theorem~3.2 of ETH-16-a, promissory obligation arises from the creation of reasonable reliance, and Principle~P1 of ETH-16-a states the conditions under which such creation obligates. Where those conditions fail, the source of promissory obligation is absent. Other duties, such as the duty not to defeat expectations one has deliberately created in third parties, are not excluded, and they bind according to their own conditions.
\end{proof}

\section{Objection and reply}

\begin{objection}
A deathbed promise binds although the promisee will never be in a position to rely on it, since she will be dead when the time for performance comes. So the obligation cannot be grounded in reliance.
\end{objection}

\begin{reply}
The dying promisee wants assurance and is entitled, in virtue of the promise, to expect performance and to make what plans she can on it, for instance by forgoing other arrangements. Principle~P1 of ETH-16-a is satisfied when the promise is made, and by Principle~P2 of ETH-16-a the obligation does not lapse when the psychological states the promise produced in her cease. Whether it persists after her death raises the further question, which ETH-16-a leaves open, whether the dead can have entitlements; EXI-13 takes up that question for duties to the dead in general. The case does not show that the obligation has a source other than reasonable reliance.
\end{reply}

\section{Conclusion}

Promises without actual reliance bind, because the promisor creates an entitlement, not a state of mind. Promises that fail the conditions of the reliance principle, such as unwanted, extorted and self-addressed promises, do not bind as promises, though their makers may be bound derivatively to those who rely on them. Both verdicts follow from the account of ETH-16-a, on which actual reliance is neither necessary for the creation of promissory obligation nor for its persistence, and reasonable reliance is its source.

\begin{thebibliography}{9}\small
\bibitem[Hume(1739)]{Hume1739} Hume, D. (1739--1740). \emph{A Treatise of Human Nature}. London: John Noon. Cited by book, part and section.
\bibitem[Scanlon(1998)]{Scanlon1998} Scanlon, T.~M. (1998). \emph{What We Owe to Each Other}. Cambridge, MA: Harvard University Press.
\end{thebibliography}

\end{document}
